\subsection*{The canonical pattern: 1 producer + 1 consumer}

The most common use of \texttt{inthreads.\allowbreak{}async} is the \textbf{1P1C} (one producer, one consumer) split for matmul:

\begin{lstlisting}[language=croqtile,caption={The canonical pattern: 1 producer 1 consumer}]
parallel p1 by 2 : group-4 {
  inthreads.async (p1 == 0) {
  }
  inthreads.async (p1 == 1) {
  }
}
\end{lstlisting}

\textbf{\texttt{parallel p1 by 2 :\allowbreak{} group-\allowbreak{}4}}

\textbf{\texttt{inthreads.\allowbreak{}async (p1 == 0)}}

\textbf{\texttt{inthreads.\allowbreak{}async (p1 == 1)}} --- The two blocks are separate programs sharing an address space.

But sharing an address space is exactly what makes this dangerous. Without coordination, the consumer might read a buffer before the producer has finished writing it. This is where events come in.
